\documentclass[10pt,twoside,openright]{report}
\usepackage{anyfontsize}
\AtBeginDocument{\fontsize{8.5}{10.6}\selectfont}

% ===================== ENGINE GUARD =====================
\usepackage{iftex}
\ifXeTeX\else
  \errmessage{This document must be compiled with XeLaTeX.}
\fi

% ===================== PAGE LAYOUT =====================
\usepackage[top=2cm,bottom=2.5cm,left=3cm,right=2cm]{geometry}
\setlength{\headheight}{15pt}

% ===================== CORE PACKAGES =====================
\usepackage{graphicx}
\usepackage{float}
\usepackage{array}
\usepackage{tabularx}
\usepackage{booktabs}
\usepackage{multirow}
\usepackage{threeparttable}
\usepackage{colortbl}
\usepackage{enumitem}
\usepackage{amsmath}
\usepackage{amssymb}
\usepackage{ragged2e}
\usepackage{longtable}
\usepackage{pdflscape}
\usepackage{xparse}
\usepackage{etoolbox}

% ===================== VISUAL DESIGN =====================
\usepackage[table]{xcolor}
\usepackage{tcolorbox}
\tcbuselibrary{breakable,skins,theorems}
\usepackage{framed}
\usepackage{titlesec}
\usepackage{tocloft}

% ===================== TIKZ FOR DIAGRAMS AND NAVIGATION =====================
\usepackage{tikz}
\usetikzlibrary{shapes,arrows,positioning,calc,decorations.pathreplacing,patterns,fit,backgrounds,shapes.geometric}

% ===================== HYPERREF FOR NAVIGATION =====================
\usepackage{hyperref}
\hypersetup{
  colorlinks=true,
  linkcolor=blue,
  filecolor=magenta,
  urlcolor=cyan,
  pdftitle={Ghost Autonomy Technical Documentation},
  pdfauthor={Ghost Autonomy},
  pdfsubject={EPU Architecture and PICAPD ISA},
  bookmarksopen=true,
  bookmarksdepth=3
}

% ===================== LISTINGS FOR CODE =====================
\usepackage{listings}
\lstset{
  basicstyle=\ttfamily\footnotesize,
  breaklines=true,
  frame=single,
  numbers=left,
  numberstyle=\tiny,
  showstringspaces=false
}

% ===================== FANCYHDR (must load BEFORE xepersian) =====================
\usepackage{fancyhdr}

% ===================== PERSIAN (XeLaTeX) — MUST BE LAST PACKAGE =====================
\usepackage{xepersian}
\setRTL
\TeXXeTstate=1

% ============ FONT SETUP ============
\setlatintextfont[
  Scale=0.95,
  Numbers=Lining,
  Ligatures=TeX
]{TeX Gyre Termes}

\settextfont[
  Path=fonts/,
  Extension=.ttf,
  UprightFont=*,
  BoldFont=*Bd,
  ItalicFont=*It,
  BoldItalicFont=*BdIt,
  Script=Arabic,
  Language=Persian
]{XB-Yagut}

\setdigitfont[
  Path=fonts/,
  Extension=.ttf,
  UprightFont=*,
  BoldFont=*Bd,
  ItalicFont=*It,
  BoldItalicFont=*BdIt,
  Numbers=Proportional,
  Script=Arabic,
  Language=Persian
]{XB-Yagut}
\setmathdigitfont[
  Path=fonts/,
  Extension=.ttf,
  UprightFont=*,
  BoldFont=*,
  ItalicFont=*,
  BoldItalicFont=*,
  Script=Arabic,
  Language=Persian
]{XB-Khoramshahr}

\providecommand{\DefaultMathsDigits}{}
\DefaultMathsDigits

% ===================== GHOST PROJECT COLOR SCHEME (8 PARTS) =====================
% Part I: Executive & Strategic
\definecolor{partExecutive}{RGB}{44,62,80}
\definecolor{chapterStrategy}{RGB}{70,90,120}
\definecolor{sectionMarket}{RGB}{90,120,150}
% Part II: Philosophy of Design (NEW)
\definecolor{partPhilosophy}{RGB}{20,90,120}
\definecolor{chapterPhilosophy}{RGB}{40,120,150}
\definecolor{sectionConcept}{RGB}{70,150,180}
% Part III: Strategic Assets & Global Partnerships (NEW)
\definecolor{partStrategic}{RGB}{120,70,50}
\definecolor{chapterStrategic}{RGB}{150,90,60}
\definecolor{sectionStrategic}{RGB}{180,120,80}
% Part IV: Technical Foundation
\definecolor{partTechnical}{RGB}{34,80,64}
\definecolor{chapterMath}{RGB}{60,110,90}
\definecolor{sectionPhysics}{RGB}{80,140,110}
% Part V: Architecture
\definecolor{partArchitecture}{RGB}{80,50,90}
\definecolor{chapterHardware}{RGB}{110,70,120}
\definecolor{sectionISA}{RGB}{140,90,150}
% Part VI: Financial & Business
\definecolor{partFinancial}{RGB}{150,80,40}
\definecolor{chapterBusiness}{RGB}{180,110,60}
\definecolor{sectionEconomics}{RGB}{210,140,80}
% Part VII: Validation & Results
\definecolor{partValidation}{RGB}{40,110,110}
\definecolor{chapterResults}{RGB}{60,140,140}
\definecolor{sectionBenchmark}{RGB}{80,170,170}
% Part VIII: Appendices & References
\definecolor{partAppendix}{RGB}{100,100,100}
\definecolor{chapterReference}{RGB}{120,120,120}
\definecolor{sectionGlossary}{RGB}{140,140,140}
% Shared utility colours
\definecolor{headerColor}{RGB}{44,62,80}
\definecolor{accentColor}{RGB}{90,70,70}
\definecolor{warningRed}{RGB}{180,40,40}
\definecolor{successGreen}{RGB}{40,150,80}
\definecolor{criticalOrange}{RGB}{200,100,30}
\definecolor{conceptBox}{RGB}{240,244,248}
\definecolor{warningBox}{RGB}{250,246,240}
\definecolor{exampleBox}{RGB}{242,248,244}
\definecolor{summaryBox}{RGB}{248,246,242}
\definecolor{processBox}{RGB}{242,248,250}
\definecolor{requirementBox}{RGB}{245,245,250}
\definecolor{criticalBox}{RGB}{255,240,240}
\definecolor{successBox}{RGB}{240,255,245}
\definecolor{tableHeader}{RGB}{44,62,80}
\definecolor{tableRow1}{RGB}{255,255,255}
\definecolor{tableRow2}{RGB}{248,249,250}

% ===================== BIDIRECTIONAL TEXT FIXES =====================
\makeatletter
\DeclareRobustCommand{\lat}[1]{%
  \ifHy@pdfstring
    #1%
  \else
    \beginL\setlatin #1\endL%
  \fi
}
\makeatother

\DeclareRobustCommand{\faen}[2]{%
  #1\ / \lat{#2}%
}

\newcommand{\english}[1]{\lat{#1}}
\newcommand{\latinbold}[1]{\lat{\textbf{#1}}}
\newcommand{\latinnum}[1]{\lat{#1}}
\newcommand{\USD}[1]{\lat{\$#1}}
\newcommand{\percent}[1]{\lat{#1\%}}
\newcommand{\technical}[1]{\lat{\textsc{#1}}}
\newcommand{\instruction}[1]{\lat{\texttt{#1}}}
\newcommand{\opcode}[1]{\lat{\texttt{#1}}}

\newenvironment{englishpar}{%
  \par\begingroup%
  \beginL%
  \setlatin%
}{%
  \endL%
  \endgroup%
  \par%
}


% ===================== NAVIGATION SYSTEM =====================
% Dynamic TOC button (context-aware return to a specific TOC entry)
\newcommand{\dynamictocbutton}[1]{%
  \begin{tikzpicture}[remember picture, overlay]
    \node[anchor=north east, xshift=-0.3cm, yshift=-1cm] at (current page.north east) {%
      \hyperlink{toc:#1}{%
        \begin{tikzpicture}
          \node[
            draw=headerColor,
            fill=conceptBox,
            line width=1pt,
            rounded corners=3pt,
            minimum width=0.8cm,
            minimum height=0.8cm,
            align=center,
            font=\tiny\bfseries,
            text=headerColor
          ] {\lat{TOC}};
        \end{tikzpicture}%
      }%
    };
  \end{tikzpicture}%
}

% Default TOC button (main hub)
\newcommand{\tocbutton}{\dynamictocbutton{maintoc}}

% ===================== 8-NODE HEXAGON NAVIGATION MAP =====================
\newcommand{\navmap}[1]{%
\begin{center}
\begin{tikzpicture}[
    hex/.style={regular polygon, regular polygon sides=6, minimum size=2cm, draw=black, line width=1pt},
    hexactive/.style={regular polygon, regular polygon sides=6, minimum size=2cm, draw=red, line width=2pt, fill=yellow!20},
    label/.style={font=\tiny\bfseries}
]
    \coordinate (center) at (0,0);
    % 8-node layout: top row (4) + bottom row (4)
    \coordinate (exec) at (-2.4, 1.2);
    \coordinate (phil) at (-0.8, 1.2);
    \coordinate (strat) at (0.8, 1.2);
    \coordinate (tech) at (2.4, 1.2);
    \coordinate (arch) at (-2.4, -1.2);
    \coordinate (fin) at (-0.8, -1.2);
    \coordinate (val) at (0.8, -1.2);
    \coordinate (app) at (2.4, -1.2);

    \ifstrequal{#1}{executive}{%
        \node[hexactive, label] (hexExec) at (exec) {\hyperlink{part:executive}{\lat{I}}};
    }{%
        \node[hex, fill=partExecutive!10, label] (hexExec) at (exec) {\hyperlink{part:executive}{\lat{I}}};
    }

    \ifstrequal{#1}{philosophy}{%
        \node[hexactive, label] (hexPhil) at (phil) {\hyperlink{part:philosophy}{\lat{II}}};
    }{%
        \node[hex, fill=partPhilosophy!10, label] (hexPhil) at (phil) {\hyperlink{part:philosophy}{\lat{II}}};
    }

    \ifstrequal{#1}{strategic}{%
        \node[hexactive, label] (hexStrat) at (strat) {\hyperlink{part:strategic}{\lat{III}}};
    }{%
        \node[hex, fill=partStrategic!10, label] (hexStrat) at (strat) {\hyperlink{part:strategic}{\lat{III}}};
    }

    \ifstrequal{#1}{technical}{%
        \node[hexactive, label] (hexTech) at (tech) {\hyperlink{part:technical}{\lat{IV}}};
    }{%
        \node[hex, fill=partTechnical!10, label] (hexTech) at (tech) {\hyperlink{part:technical}{\lat{IV}}};
    }

    \ifstrequal{#1}{architecture}{%
        \node[hexactive, label] (hexArch) at (arch) {\hyperlink{part:architecture}{\lat{V}}};
    }{%
        \node[hex, fill=partArchitecture!10, label] (hexArch) at (arch) {\hyperlink{part:architecture}{\lat{V}}};
    }

    \ifstrequal{#1}{financial}{%
        \node[hexactive, label] (hexFin) at (fin) {\hyperlink{part:financial}{\lat{VI}}};
    }{%
        \node[hex, fill=partFinancial!10, label] (hexFin) at (fin) {\hyperlink{part:financial}{\lat{VI}}};
    }

    \ifstrequal{#1}{validation}{%
        \node[hexactive, label] (hexVal) at (val) {\hyperlink{part:validation}{\lat{VII}}};
    }{%
        \node[hex, fill=partValidation!10, label] (hexVal) at (val) {\hyperlink{part:validation}{\lat{VII}}};
    }

    \ifstrequal{#1}{appendices}{%
        \node[hexactive, label] (hexApp) at (app) {\hyperlink{part:appendices}{\lat{VIII}}};
    }{%
        \node[hex, fill=partAppendix!10, label] (hexApp) at (app) {\hyperlink{part:appendices}{\lat{VIII}}};
    }

    \node[circle, draw=headerColor, fill=white, line width=1.5pt, minimum size=1.8cm, font=\scriptsize\bfseries] at (center) {%
      \hyperlink{maintoc}{%
        \begin{tabular}{@{}c@{}}%
          \lat{Ghost}\\[2pt]%
          \lat{Autonomy}%
        \end{tabular}%
      }%
    };

    \node[below=0.1cm of hexExec, font=\tiny] {\color{partExecutive}\lat{Executive}};
    \node[below=0.1cm of hexPhil, font=\tiny] {\color{partPhilosophy}\lat{Philosophy}};
    \node[below=0.1cm of hexStrat, font=\tiny] {\color{partStrategic}\lat{Strategic}};
    \node[below=0.1cm of hexTech, font=\tiny] {\color{partTechnical}\lat{Technical}};
    \node[below=0.1cm of hexArch, font=\tiny] {\color{partArchitecture}\lat{Architecture}};
    \node[below=0.1cm of hexFin, font=\tiny] {\color{partFinancial}\lat{Financial}};
    \node[below=0.1cm of hexVal, font=\tiny] {\color{partValidation}\lat{Validation}};
    \node[below=0.1cm of hexApp, font=\tiny] {\color{partAppendix}\lat{Appendices}};
\end{tikzpicture}
\end{center}
}

% Navigation widget combining TOC button and hexagon map
\newcommand{\navwidget}[2]{%
  \dynamictocbutton{#1}%  % TOC return button (to exact entry)
  \vspace{0.5cm}%         % Spacing
  \navmap{#2}%            % Hexagon map
}


% ===================== HEADER/FOOTER WITH NAVIGATION =====================
\pagestyle{fancy}
\fancyhf{}
\fancyhead[RE]{\small\nouppercase{\rightmark}}
\fancyhead[LO]{\small\nouppercase{\leftmark}}
\fancyhead[LE,RO]{\bfseries\thepage}
\fancyfoot[C]{\footnotesize\color{gray}\lat{Ghost Autonomy — Complete Technical Documentation v2.0}}
\renewcommand{\headrulewidth}{0.4pt}
\renewcommand{\footrulewidth}{0.3pt}

\fancypagestyle{plain}{%
  \fancyhf{}
  \fancyfoot[C]{\bfseries\thepage}
  \renewcommand{\headrulewidth}{0pt}
  \renewcommand{\footrulewidth}{0pt}
}

% ===================== CHAPTER/SECTION FORMATTING =====================
\titleformat{\part}[display]
  {\normalfont\Huge\bfseries\color{headerColor}\centering}
  {\partname\ \thepart}
  {20pt}
  {\Huge}
\titlespacing*{\part}{0pt}{-20pt}{40pt}

\titleformat{\chapter}[display]
  {\normalfont\huge\bfseries}
  {\flushright\chaptertitlename\ \thechapter}
  {20pt}
  {\Huge\flushright\color{headerColor}}
\titlespacing*{\chapter}{0pt}{-20pt}{30pt}

\titleformat{\section}
  {\normalfont\Large\bfseries\color{headerColor}}
  {\thesection}{1em}{}

\titleformat{\subsection}
  {\normalfont\large\bfseries\color{chapterStrategy}}
  {\thesubsection}{1em}{}

\titleformat{\subsubsection}
  {\normalfont\normalsize\bfseries\color{sectionMarket}}
  {\thesubsubsection}{1em}{}

% ===================== ENHANCED TABLE OF CONTENTS =====================
\setcounter{tocdepth}{3}
\setcounter{secnumdepth}{3}

\renewcommand{\cftpartfont}{\bfseries\LARGE\color{headerColor}}
\renewcommand{\cftpartpagefont}{\bfseries\LARGE\color{headerColor}}
\setlength{\cftbeforepartskip}{25pt}
\setlength{\cftpartindent}{0pt}

\renewcommand{\cftchapfont}{\bfseries\large\color{chapterStrategy}}
\renewcommand{\cftchappagefont}{\bfseries\large\color{chapterStrategy}}
\setlength{\cftbeforechapskip}{10pt}
\setlength{\cftchapindent}{15pt}

\renewcommand{\cftsecfont}{\normalsize\color{sectionMarket}}
\renewcommand{\cftsecpagefont}{\normalsize\color{sectionMarket}}
\setlength{\cftbeforesecskip}{4pt}
\setlength{\cftsecindent}{30pt}

\renewcommand{\cftsubsecfont}{\small\color{gray}}
\renewcommand{\cftsubsecpagefont}{\small\color{gray}}
\setlength{\cftsubsecindent}{45pt}

\renewcommand{\cftchapleader}{\cftdotfill{\cftdotsep}}
\renewcommand{\cftsecleader}{\cftdotfill{\cftdotsep}}

\renewcommand{\contentsname}{\hfill\Huge\bfseries\color{headerColor}


فهرست مطالب جامع\\[10pt]
\large\color{chapterStrategy}\lat{Complete Table of Contents}\hfill}

% ===================== AUTOMATIC TOC TARGET INJECTION =====================
\makeatletter
\let\oldaddcontentsline\addcontentsline
\renewcommand{\addcontentsline}[3]{%
  \ifstrequal{#1}{toc}{%
    \begingroup
      \let\GA@toctarget\@empty
      \ifstrequal{#2}{part}{\edef\GA@toctarget{toc:part\thepart}}{}%
      \ifstrequal{#2}{chapter}{\edef\GA@toctarget{toc:chap\thechapter}}{}%
      \ifstrequal{#2}{section}{\edef\GA@toctarget{toc:sec\thesection}}{}%
      \ifstrequal{#2}{subsection}{\edef\GA@toctarget{toc:subsec\thesubsection}}{}%
      \ifstrequal{#2}{subsubsection}{\edef\GA@toctarget{toc:subsubsec\thesubsubsection}}{}%
      \ifx\GA@toctarget\@empty\else
        \addtocontents{#1}{\protect\hypertarget{\GA@toctarget}{}}%
      \fi
    \endgroup
  }{}%
  \oldaddcontentsline{#1}{#2}{#3}%
}
\makeatother

% ===================== RETURN TO TOC WITH NAVIGATION =====================
\newcommand{\gototoc}[2]{%
\vspace{1em}
\noindent\begin{center}
\navwidget{#1}{#2}
\end{center}
\vspace{1em}
}

% ===================== PART SEPARATOR =====================
\newcommand{\partseparator}[2]{%
\clearpage
\thispagestyle{empty}
\vspace*{\fill}
\begin{center}
{\Huge\bfseries\color{#1}═══════════════════════════════════}\\[20pt]
{\LARGE\bfseries بخش جدید آغاز می‌شود\\[5pt]\lat{New Part Begins}}\\[20pt]
{\Huge\bfseries\color{#1}═══════════════════════════════════}\\[30pt]
\navmap{#2}
\end{center}
\vspace*{\fill}
\clearpage
}

% ===================== PAGE ALLOCATION =====================
\newcounter{allocpage}
\newif\iffastcompile
\fastcompiletrue % ← SET TO \fastcompilefalse FOR FULL 3,140-PAGE BUILD
\newcount\allocTargetCount

\newcommand{\allocatepages}[5]{%
  \allocTargetCount=#1\relax
  \iffastcompile
    \ifnum\allocTargetCount>1 \allocTargetCount=1\relax\fi
  \fi
  \setcounter{allocpage}{0}%
  \loop\ifnum\value{allocpage}<\the\allocTargetCount
    \stepcounter{allocpage}%
    \clearpage
    \thispagestyle{fancy}
    \vspace*{2cm}
    \begin{center}
    \begin{tcolorbox}[
      colback=#3!5,
      colframe=#3,
      width=0.9\textwidth,
      arc=3mm,
      boxrule=1.5pt,
      title={\large\bfseries #2}
    ]
    \vspace{1cm}
    \centering
    {\large\bfseries صفحه \arabic{allocpage}\ از #1\\[5pt]\lat{Page \arabic{allocpage}\ of #1}}
    \vspace{1cm}
    \end{tcolorbox}
    \end{center}
    \vspace{2cm}
    \gototoc{#4}{#5}
  \repeat
}

% ===================== DOCUMENT STRUCTURE =====================
\begin{document}

%═══════════════════════════════════════════════════════════════════════════
% TITLE PAGE
%═══════════════════════════════════════════════════════════════════════════
\begin{titlepage}
\centering
\vspace*{2cm}

{\Huge\bfseries\color{partExecutive}\lat{Ghost Autonomy}\par}
\vspace{0.3cm}
{\LARGE\bfseries\color{headerColor}سیستم خودمختار رانندگی مبتنی بر معماری EPU\par}
\vspace{0.3cm}
{\Large\color{chapterStrategy}\lat{Autonomous Driving System Based on EPU Architecture}\par}

\vspace{3cm}

\begin{tcolorbox}[
  colback=conceptBox,
  colframe=partExecutive,
  width=0.8\textwidth,
  arc=3mm,
  boxrule=1pt
]
\centering
\large
\textbf{دکتر محسن دیرباز}\\[5pt]
\lat{Founder \& Chief Architect \& Didas International Group - VP Semiconductor}\\[10pt]
\normalsize
با همکاری تیم فنی \lat{Ghost Autonomy}\\[5pt]
\lat{Technical Documentation v2.0}\\[15pt]
{\footnotesize\color{gray}
\lat{February 2026}\\
تحت پوشش حقوقی \lat{Patent Pending} و \lat{Trade Secret}
}
\end{tcolorbox}

\vfill

\begin{tikzpicture}[scale=1.5]
  \draw[fill=conceptBox, draw=headerColor, line width=2pt] (0,0) circle (2.5cm);
  \node at (0,0.8) {\Large\bfseries\lat{Ghost}};
  \node at (0,0.2) {\Large\bfseries\lat{Autonomy}};
  \node at (0,-0.5) {\large\lat{EPU}};
  \node at (0,-1) {\footnotesize\lat{Architecture}};
\end{tikzpicture}

\vspace{2cm}

{\LARGE\bfseries\color{chapterMath}مستندات فنی جامع\par}
\vspace{0.2cm}
{\Large\color{sectionPhysics}\lat{Complete Technical Documentation}\par}

\vspace{1cm}

{\large\bfseries
نسخه ۲٫۰ — معماری PICAPD و سیستم EPU\\
\lat{Version 2.0 — PICAPD Architecture \& EPU System}
\par}

\vfill

{\large\color{gray}
\lat{© 2026 Ghost Autonomy. All Rights Reserved.}\\[5pt]
تاریخ: فوریه ۲۰۲۶ / \lat{Date: February 2026}
\par}
\end{titlepage}

%═══════════════════════════════════════════════════════════════════════════
% LEGAL PAGE
%═══════════════════════════════════════════════════════════════════════════
\clearpage
\thispagestyle{empty}
\vspace*{3cm}

\begin{center}
{\Large\bfseries\color{headerColor}اطلاعیه حقوقی / \lat{Legal Notice}\par}
\end{center}

\vspace{1cm}

\begin{tcolorbox}[
  colback=warningBox,
  colframe=warningRed,
  arc=3mm,
  boxrule=1pt
]
\begin{englishpar}
\textbf{Copyright \lat{© 2026 Ghost Autonomy.}} All rights reserved.

This document contains proprietary and confidential information. Unauthorized reproduction, distribution, or disclosure is strictly prohibited.

\vspace{0.5cm}

\textbf{Patent Status:} Patents pending on PICAPD ISA, EPU architecture, Byzantine consensus protocols, and moment-based population dynamics.

\vspace{0.5cm}

\textbf{Trademarks:} Ghost Autonomy, PICAPD, EPU, and related marks are trademarks of Ghost Autonomy.

\vspace{0.5cm}

\textbf{Export Control:} This document may contain technical data subject to export control regulations.
\end{englishpar}
\end{tcolorbox}

\vfill

%═══════════════════════════════════════════════════════════════════════════
% TABLE OF CONTENTS WITH 8-PART NAVIGATION MAP
%═══════════════════════════════════════════════════════════════════════════
\clearpage
\hypertarget{maintoc}{}
\hypertarget{toc:maintoc}{}
\pdfbookmark[0]{\lat{Table of Contents}}{maintoc}

\begin{center}
\vspace*{1cm}
{\Huge\bfseries\color{headerColor}فهرست مطالب جامع\par}
\vspace{0.3cm}
{\Large\color{chapterStrategy}\lat{Complete Table of Contents}\par}
\vspace{1cm}

\navmap{none}

\vspace{1cm}

\begin{tcolorbox}[
  colback=conceptBox,
  colframe=headerColor,
  width=0.95\textwidth,
  arc=5mm,
  boxrule=2pt,
  title={\Large\bfseries راهنمای رنگ‌بندی / \lat{Color Legend}}
]
\begin{tabular}{ll}
\colorbox{partExecutive!20}{\strut\hspace{1cm}} & \textbf{بخش اول: اجرایی و استراتژیک / \lat{Part I: Executive \& Strategic}} \\[5pt]
\colorbox{partPhilosophy!20}{\strut\hspace{1cm}} & \textbf{بخش دوم: فلسفه طراحی / \lat{Part II: Philosophy of Design}} \\[5pt]
\colorbox{partStrategic!20}{\strut\hspace{1cm}} & \textbf{بخش سوم: دارایی‌های استراتژیک / \lat{Part III: Strategic Assets}} \\[5pt]
\colorbox{partTechnical!20}{\strut\hspace{1cm}} & \textbf{بخش چهارم: مبانی فنی / \lat{Part IV: Technical Foundation}} \\[5pt]
\colorbox{partArchitecture!20}{\strut\hspace{1cm}} & \textbf{بخش پنجم: معماری / \lat{Part V: Architecture}} \\[5pt]
\colorbox{partFinancial!20}{\strut\hspace{1cm}} & \textbf{بخش ششم: مالی / \lat{Part VI: Financial}} \\[5pt]
\colorbox{partValidation!20}{\strut\hspace{1cm}} & \textbf{بخش هفتم: اعتبارسنجی / \lat{Part VII: Validation}} \\[5pt]
\colorbox{partAppendix!20}{\strut\hspace{1cm}} & \textbf{بخش هشتم: پیوست‌ها / \lat{Part VIII: Appendices}}
\end{tabular}
\end{tcolorbox}
\end{center}

\vspace{0.5cm}

\tableofcontents

\clearpage

%═══════════════════════════════════════════════════════════════════════════
% PART I: EXECUTIVE & STRATEGIC (250 pages) — Chapters 1–7
%═══════════════════════════════════════════════════════════════════════════
\partseparator{partExecutive}{executive}

\part{\faen{اجرایی و استراتژیک}{Executive \& Strategic}}
\label{part:executive}
\hypertarget{part:executive}{}

\chapter{\faen{خلاصه اجرایی}{Executive Summary}}
\label{chap:executive_summary}

\section{\faen{چشم‌انداز}{Vision}}
\allocatepages{10}{\faen{فصل ۱، بخش ۱: چشم‌انداز}{Ch.\ 1, Sec.\ 1: Vision}}{chapterStrategy}{sec1.1}{executive}

\section{\faen{اهداف کلان}{Strategic Goals}}
\allocatepages{10}{\faen{فصل ۱، بخش ۲: اهداف}{Ch.\ 1, Sec.\ 2: Goals}}{chapterStrategy}{sec1.2}{executive}

\section{\faen{پیشنهاد ارزش}{Value Proposition}}
\allocatepages{10}{\faen{فصل ۱، بخش ۳: پیشنهاد ارزش}{Ch.\ 1, Sec.\ 3: Value Prop}}{chapterStrategy}{sec1.3}{executive}

\section{\faen{خلاصه فنی}{Technical Summary}}
\allocatepages{10}{\faen{فصل ۱، بخش ۴: خلاصه فنی}{Ch.\ 1, Sec.\ 4: Tech Summary}}{chapterStrategy}{sec1.4}{executive}

\chapter{\faen{تحلیل بازار}{Market Analysis}}
\label{chap:market_analysis}

\section{\faen{اندازه بازار}{Market Size}}
\allocatepages{10}{\faen{فصل ۲، بخش ۱: اندازه بازار}{Ch.\ 2, Sec.\ 1: Market Size}}{sectionMarket}{sec2.1}{executive}

\section{\faen{رقبا}{Competitive Landscape}}
\allocatepages{15}{\faen{فصل ۲، بخش ۲: رقبا}{Ch.\ 2, Sec.\ 2: Competitors}}{sectionMarket}{sec2.2}{executive}

\section{\faen{روندها}{Trends}}
\allocatepages{10}{\faen{فصل ۲، بخش ۳: روندها}{Ch.\ 2, Sec.\ 3: Trends}}{sectionMarket}{sec2.3}{executive}

\section{\faen{فرصت‌ها}{Opportunities}}
\allocatepages{5}{\faen{فصل ۲، بخش ۴: فرصت‌ها}{Ch.\ 2, Sec.\ 4: Opportunities}}{sectionMarket}{sec2.4}{executive}

\chapter{\faen{استراتژی محصول}{Product Strategy}}
\label{chap:product_strategy}
\allocatepages{40}{\faen{فصل ۳: استراتژی محصول}{Chapter 3: Product Strategy}}{chapterStrategy}{chap3}{executive}

\chapter{\faen{مزیت رقابتی}{Competitive Advantage}}
\label{chap:competitive_advantage}
\allocatepages{30}{\faen{فصل ۴: مزیت رقابتی}{Chapter 4: Competitive Advantage}}{chapterStrategy}{chap4}{executive}

\chapter{\faen{نقشه راه}{Roadmap}}
\label{chap:roadmap}
\allocatepages{40}{\faen{فصل ۵: نقشه راه}{Chapter 5: Roadmap}}{chapterStrategy}{chap5}{executive}

\chapter{\faen{تیم}{Team}}
\label{chap:team}
\allocatepages{30}{\faen{فصل ۶: تیم}{Chapter 6: Team}}{chapterStrategy}{chap6}{executive}

\chapter{\faen{مدیریت ریسک}{Risk Management}}
\label{chap:risk_management}
\allocatepages{30}{\faen{فصل ۷: مدیریت ریسک}{Chapter 7: Risk Management}}{chapterStrategy}{chap7}{executive}

%═══════════════════════════════════════════════════════════════════════════
% PART II: PHILOSOPHY OF DESIGN AND CONCEPTUALIZATION (190 pages) — Chapter 8
%═══════════════════════════════════════════════════════════════════════════
\partseparator{partPhilosophy}{philosophy}

\part{\faen{فلسفه طراحی و مفهوم‌سازی}{Philosophy of Design and Conceptualization}}
\label{part:philosophy}
\hypertarget{part:philosophy}{}

\chapter{\faen{فلسفه طراحی و مفهوم‌سازی}{Philosophy of Design and Conceptualization}}
\label{chap:philosophy_design}

\section{\faen{انگیزه، دامنه و قرارداد خواننده}{Motivation, Scope, and the Reader Contract}}
\allocatepages{25}{\faen{فصل ۸، بخش ۱: انگیزه و دامنه}{Ch.\ 8, Sec.\ 1: Motivation \& Scope}}{chapterPhilosophy}{sec8.1}{philosophy}

\subsection{\faen{چرا «فلسفه طراحی» پیش‌نیاز معماری است}{Why Design Philosophy Precedes Architecture}}
\subsection{\faen{مرزهای دامنه و «غیرهدف‌ها»}{Scope Boundaries and Non-Goals}}
\subsection{\faen{قرارداد خواننده: چگونه این سند را بخوانیم}{Reader Contract: How to Read This Document}}
\subsection{\faen{هم‌ترازی با نقشه پیاده‌سازی}{Alignment with the Implementation Map}}

\section{\faen{مدل مفهومی سامانه: عامل، جهان و علیت}{Conceptual System Model: Agent, World, and Causality}}
\allocatepages{30}{\faen{فصل ۸، بخش ۲: مدل مفهومی}{Ch.\ 8, Sec.\ 2: Conceptual Model}}{chapterPhilosophy}{sec8.2}{philosophy}

\subsection{\faen{تعریف سامانه خودمختار به‌عنوان مسئله کنترل}{Autonomy as a Control Problem}}
\subsection{\faen{مرز سامانه و محیط: قراردادهای مشاهده و عمل}{System Boundary: Observation/Action Contracts}}
\subsection{\faen{زمان، حالت و تاریخچه}{Time, State, and History}}
\subsection{\faen{علیت، پیش‌بینی و سازگاری}{Causality, Prediction, and Consistency}}

\section{\faen{اصول طراحی: سازگاری، سادگی و مقیاس‌پذیری}{Design Principles: Coherence, Simplicity, and Scalability}}
\allocatepages{35}{\faen{فصل ۸، بخش ۳: اصول طراحی}{Ch.\ 8, Sec.\ 3: Design Principles}}{chapterPhilosophy}{sec8.3}{philosophy}

\subsection{\faen{تفکیک نگرانی‌ها و مرزبندی مسئولیت‌ها}{Separation of Concerns and Responsibility Boundaries}}
\subsection{\faen{قرارداد-محور: تعریف دقیق ورودی/خروجی}{Contract-First Interfaces}}
\subsection{\faen{قطعیّت و بازتولیدپذیری}{Determinism and Reproducibility}}
\subsection{\faen{قابلیت توضیح و قابلیت بازرسی}{Explainability and Inspectability}}
\subsection{\faen{پایداری کارایی و قابل‌حمل بودن عملکرد}{Performance Portability}}

\section{\faen{لایه‌های انتزاع و قراردادهای اطلاعاتی}{Abstraction Layers and Information Contracts}}
\allocatepages{30}{\faen{فصل ۸، بخش ۴: لایه‌های انتزاع}{Ch.\ 8, Sec.\ 4: Abstraction Layers}}{chapterPhilosophy}{sec8.4}{philosophy}

\subsection{\faen{از معنا تا سیگنال: سلسله‌مراتب نمایش}{From Semantics to Signals: Representation Hierarchy}}
\subsection{\faen{قراردادهای میانی بین ادراک، پیش‌بینی و تصمیم}{Mid-Layer Contracts: Perception--Prediction--Decision}}
\subsection{\faen{قرارداد ISA: چرا PICAPD نقطه ثقل است}{The ISA Contract: Why PICAPD is the Fulcrum}}
\subsection{\faen{قرارداد زمان‌بندی و حافظه: هزینه و صحت}{Timing/Memory Contracts: Cost vs Correctness}}

\section{\faen{عدم‌قطعیت، ریسک و پرونده ایمنی}{Uncertainty, Risk, and the Safety Case}}
\allocatepages{30}{\faen{فصل ۸، بخش ۵: ریسک و ایمنی}{Ch.\ 8, Sec.\ 5: Risk \& Safety}}{chapterPhilosophy}{sec8.5}{philosophy}

\subsection{\faen{تفکیک عدم‌قطعیت معرفتی و تصادفی}{Epistemic vs Aleatoric Uncertainty}}
\subsection{\faen{بودجه ریسک و سیاست‌های شکستِ امن}{Risk Budgets and Fail-Safe Policies}}
\subsection{\faen{قابلیت اثبات‌پذیری: از شواهد تا ادعا}{Assurance: From Evidence to Claims}}
\subsection{\faen{تعامل انسان و سامانه در حلقه تصمیم}{Human-in-the-Loop Boundaries}}

\section{\faen{مستندسازی به‌عنوان معماری}{Documentation as Architecture}}
\allocatepages{25}{\faen{فصل ۸، بخش ۶: مستندسازی}{Ch.\ 8, Sec.\ 6: Documentation}}{chapterPhilosophy}{sec8.6}{philosophy}

\subsection{\faen{ماژولار بودن سند و قراردادهای ارجاع}{Document Modularity and Reference Contracts}}
\subsection{\faen{نسخه‌بندی معناشناختی و سازگاری عقب‌رو}{Semantic Versioning and Backward Compatibility}}
\subsection{\faen{قابلیت رهگیری: پیوند نیازمندی→طراحی→پیاده‌سازی}{Traceability: Requirements to Design to Implementation}}
\subsection{\faen{حلقه بازخورد: بازبینی، آزمون و به‌روزرسانی}{Feedback Loop: Review, Test, Update}}

\section{\faen{پل مفهومی به مبانی فنی}{Conceptual Bridge to the Technical Foundations}}
\allocatepages{15}{\faen{فصل ۸، بخش ۷: پل مفهومی}{Ch.\ 8, Sec.\ 7: Bridge Forward}}{chapterPhilosophy}{sec8.7}{philosophy}

\subsection{\faen{چرا مبانی فیزیک‌الهام‌گرفته در ادامه می‌آید}{Why Physics-Inspired Foundations Come Next}}
\subsection{\faen{نگاشت اصول به ماژول‌های فنی}{Mapping Principles to Technical Modules}}
\subsection{\faen{راهنمای خواندن بخش فنی و معماری}{Reader Guide for Technical and Architecture Parts}}

%═══════════════════════════════════════════════════════════════════════════
% PART III: STRATEGIC ASSETS & GLOBAL PARTNERSHIPS (240 pages) — Chapters 9–13
%═══════════════════════════════════════════════════════════════════════════
\partseparator{partStrategic}{strategic}

\part{\faen{دارایی‌های استراتژیک و مشارکت‌های جهانی}{Strategic Assets and Global Partnerships}}
\label{part:strategic}
\hypertarget{part:strategic}{}

\chapter{\faen{پیشینه علمی: از IIT تا سیلیکون}{Scientific Provenance: From IIT to Silicon}}
\label{chap:scientific_provenance}

\section{\faen{خاستگاه نگاشت دوخطی}{Origin of the Bilinear Form Mapping}}
\allocatepages{15}{\faen{فصل ۹، بخش ۱: خاستگاه}{Ch.\ 9, Sec.\ 1: Origin}}{chapterStrategic}{sec9.1}{strategic}

\section{\faen{از گازی‌سازی زیست‌توده تا محاسبات خودران}{From Biomass Gasification to Autonomous Computing}}
\allocatepages{12}{\faen{فصل ۹، بخش ۲: گازی‌سازی}{Ch.\ 9, Sec.\ 2: Gasification}}{chapterStrategic}{sec9.2}{strategic}

\section{\faen{انتقال فناوری و ثبت شرکت در آمریکا}{Technology Transfer \& US Incorporation}}
\allocatepages{8}{\faen{فصل ۹، بخش ۳: انتقال فناوری}{Ch.\ 9, Sec.\ 3: Tech Transfer}}{chapterStrategic}{sec9.3}{strategic}

\section{\faen{همکاری‌های علمی مستمر}{Continuing Research Ties}}
\allocatepages{6}{\faen{فصل ۹، بخش ۴: همکاری علمی}{Ch.\ 9, Sec.\ 4: Research Ties}}{chapterStrategic}{sec9.4}{strategic}

\section{\faen{خط لوله استعدادها: دانشگاه صنعتی شریف}{Talent Pipeline: Sharif University}}
\allocatepages{5}{\faen{فصل ۹، بخش ۵: شریف}{Ch.\ 9, Sec.\ 5: Sharif}}{chapterStrategic}{sec9.5}{strategic}

\section{\faen{مطالعه موردی: پروفسور مهدی طاهری و imec}{Case Study: Prof.\ Mahdi Taheri and imec}}
\allocatepages{4}{\faen{فصل ۹، بخش ۶: طاهری}{Ch.\ 9, Sec.\ 6: Taheri}}{chapterStrategic}{sec9.6}{strategic}

\chapter{\faen{یکپارچگی عمودی: معدن سیلیس با خلوص بالا}{Vertical Integration: The High-Purity Silica Asset}}
\label{chap:silica_mine}

\section{\faen{زنجیره تأمین جهانی سیلیس}{Global Silica Supply Chain}}
\allocatepages{8}{\faen{فصل ۱۰، بخش ۱: زنجیره تأمین}{Ch.\ 10, Sec.\ 1: Supply Chain}}{chapterStrategic}{sec10.1}{strategic}

\section{\faen{مشخصات معدن ایران}{Characterisation of the Iranian Mine}}
\allocatepages{10}{\faen{فصل ۱۰، بخش ۲: مشخصات معدن}{Ch.\ 10, Sec.\ 2: Mine Specs}}{chapterStrategic}{sec10.2}{strategic}

\section{\faen{از سیلیس تا پلی‌سیلیکون تا ویفر}{From Silica to Polysilicon to Wafers}}
\allocatepages{8}{\faen{فصل ۱۰، بخش ۳: نقشه راه فناوری}{Ch.\ 10, Sec.\ 3: Tech Roadmap}}{chapterStrategic}{sec10.3}{strategic}

\section{\faen{ارزش استراتژیک: هزینه، امنیت و ESG}{Strategic Value: Cost, Security, ESG}}
\allocatepages{6}{\faen{فصل ۱۰، بخش ۴: ارزش استراتژیک}{Ch.\ 10, Sec.\ 4: Strategic Value}}{chapterStrategic}{sec10.4}{strategic}

\section{\faen{الگوهای مشارکت}{Partnership Models}}
\allocatepages{8}{\faen{فصل ۱۰، بخش ۵: مشارکت}{Ch.\ 10, Sec.\ 5: Partnerships}}{chapterStrategic}{sec10.5}{strategic}

\chapter{\faen{ماموریت نیمه‌هادی هند و Ghost Autonomy}{India's Semiconductor Mission and Ghost Autonomy}}
\label{chap:india_partnership}

\section{\faen{مروری بر ISM 1.0 و 2.0}{Overview of ISM 1.0 and 2.0}}
\allocatepages{10}{\faen{فصل ۱۱، بخش ۱: ISM}{Ch.\ 11, Sec.\ 1: ISM}}{chapterStrategic}{sec11.1}{strategic}

\section{\faen{هم‌ترازی با قابلیت‌های Ghost}{Alignment with Ghost's Strengths}}
\allocatepages{12}{\faen{فصل ۱۱، بخش ۲: هم‌ترازی}{Ch.\ 11, Sec.\ 2: Alignment}}{chapterStrategic}{sec11.2}{strategic}

\section{\faen{چارچوب همکاری پیشنهادی}{Proposed Collaboration Framework}}
\allocatepages{18}{\faen{فصل ۱۱، بخش ۳: چارچوب}{Ch.\ 11, Sec.\ 3: Framework}}{chapterStrategic}{sec11.3}{strategic}

\section{\faen{ارزش پیشنهادی برای هند}{Value Proposition for India}}
\allocatepages{15}{\faen{فصل ۱۱، بخش ۴: ارزش پیشنهادی}{Ch.\ 11, Sec.\ 4: Value Prop}}{chapterStrategic}{sec11.4}{strategic}

\chapter{\faen{برتری نیمه‌هادی کره جنوبی و Ghost Autonomy}{South Korea's Semiconductor Supremacy and Ghost Autonomy}}
\label{chap:korea_partnership}

\section{\faen{جایگاه کره: حافظه، فاندری، HBM}{Korea's Position: Memory, Foundry, HBM for AI}}
\allocatepages{10}{\faen{فصل ۱۲، بخش ۱: جایگاه کره}{Ch.\ 12, Sec.\ 1: Korea Position}}{chapterStrategic}{sec12.1}{strategic}

\section{\faen{سه‌گانه کره-آمریکا-ایران: کنترل صادرات}{The Korea-US-Iran Trilemma: Export Controls}}
\allocatepages{10}{\faen{فصل ۱۲، بخش ۲: سه‌گانه}{Ch.\ 12, Sec.\ 2: Trilemma}}{chapterStrategic}{sec12.2}{strategic}

\section{\faen{فرصت‌های همکاری}{Collaboration Opportunities}}
\allocatepages{20}{\faen{فصل ۱۲، بخش ۳: همکاری}{Ch.\ 12, Sec.\ 3: Collaboration}}{chapterStrategic}{sec12.3}{strategic}

\section{\faen{ارزش پیشنهادی برای کره جنوبی}{Value Proposition for South Korea}}
\allocatepages{15}{\faen{فصل ۱۲، بخش ۴: ارزش پیشنهادی}{Ch.\ 12, Sec.\ 4: Value Prop}}{chapterStrategic}{sec12.4}{strategic}

\chapter{\faen{اتحاد چهارجانبه: هند-کره-آمریکا-ایران}{The Quad-Plus Alliance: India-Korea-US-Iran}}
\label{chap:quad_alliance}

\section{\faen{ترکیب نقاط قوت مکمل}{Synthesis of Complementary Strengths}}
\allocatepages{10}{\faen{فصل ۱۳، بخش ۱: ترکیب}{Ch.\ 13, Sec.\ 1: Synthesis}}{chapterStrategic}{sec13.1}{strategic}

\section{\faen{ابتکارات چندجانبه پیشنهادی}{Proposed Multi-lateral Initiatives}}
\allocatepages{15}{\faen{فصل ۱۳، بخش ۲: ابتکارات}{Ch.\ 13, Sec.\ 2: Initiatives}}{chapterStrategic}{sec13.2}{strategic}

\section{\faen{نقشه راه پروژه آزمایشی (۲۰۲۶–۲۰۲۸)}{Roadmap to a Pilot Project (2026--2028)}}
\allocatepages{8}{\faen{فصل ۱۳، بخش ۳: نقشه راه}{Ch.\ 13, Sec.\ 3: Roadmap}}{chapterStrategic}{sec13.3}{strategic}

\section{\faen{مدیریت ریسک}{Risk Mitigation}}
\allocatepages{7}{\faen{فصل ۱۳، بخش ۴: ریسک}{Ch.\ 13, Sec.\ 4: Risk}}{chapterStrategic}{sec13.4}{strategic}

%═══════════════════════════════════════════════════════════════════════════
% PART IV: TECHNICAL FOUNDATION (800 pages) — Chapters 14–22
%═══════════════════════════════════════════════════════════════════════════
\partseparator{partTechnical}{technical}

\part{\faen{مبانی فنی}{Technical Foundation}}
\label{part:technical}
\hypertarget{part:technical}{}

\chapter{\faen{محاسبات الهام‌گرفته از فیزیک}{Physics-Inspired Computing}}
\label{chap:physics_computing}

\section{\faen{اصول فیزیکی}{Physical Principles}}
\allocatepages{20}{\faen{فصل ۱۴، بخش ۱: اصول}{Ch.\ 14, Sec.\ 1: Principles}}{chapterMath}{sec14.1}{technical}

\section{\faen{معادلات حاکم}{Governing Equations}}
\allocatepages{25}{\faen{فصل ۱۴، بخش ۲: معادلات}{Ch.\ 14, Sec.\ 2: Equations}}{chapterMath}{sec14.2}{technical}

\subsection{\faen{هذلولی}{Hyperbolic}}
\subsection{\faen{سهموی}{Parabolic}}
\subsection{\faen{انتقال}{Transport}}
\subsection{\lat{ODE}}

\section{\faen{پیوندهای دوخطی}{Bilinear Coupling}}
\allocatepages{20}{\faen{فصل ۱۴، بخش ۳: پیوندها}{Ch.\ 14, Sec.\ 3: Coupling}}{chapterMath}{sec14.3}{technical}

\section{\faen{فضاهای توپولوژیک}{Topological Vector Spaces}}
\allocatepages{15}{\faen{فصل ۱۴، بخش ۴: توپولوژی}{Ch.\ 14, Sec.\ 4: Topology}}{chapterMath}{sec14.4}{technical}

\section{\faen{نگاشت دوخطی -- رساله IIT (اشتقاق کامل)}{The Bilinear Form Mapping -- IIT Dissertation (Full Derivation)}}
\allocatepages{30}{\faen{فصل ۱۴، بخش ۵: نگاشت دوخطی IIT}{Ch.\ 14, Sec.\ 5: Bilinear Mapping}}{chapterMath}{sec14.5}{technical}

\chapter{\faen{مبانی ریاضی}{Mathematical Foundations}}
\label{chap:math_foundations}

\section{\faen{آنالیز هارمونیک}{Harmonic Analysis}}
\allocatepages{25}{\faen{فصل ۱۵، بخش ۱: هارمونیک}{Ch.\ 15, Sec.\ 1: Harmonic}}{sectionPhysics}{sec15.1}{technical}

\subsection{\faen{تبدیل فوریه}{Fourier Transform}}
\subsection{\faen{انتگرال بیضوی}{Elliptic Integrals}}
\subsection{\faen{میانگین AGM}{AGM Iteration}}

\section{\faen{تحلیل تابعی}{Functional Analysis}}
\allocatepages{25}{\faen{فصل ۱۵، بخش ۲: تحلیل تابعی}{Ch.\ 15, Sec.\ 2: Functional}}{sectionPhysics}{sec15.2}{technical}

\section{\faen{نظریه توزیع}{Distribution Theory}}
\allocatepages{25}{\faen{فصل ۱۵، بخش ۳: توزیع‌ها}{Ch.\ 15, Sec.\ 3: Distributions}}{sectionPhysics}{sec15.3}{technical}

\subsection{\faen{فضای شوارتز}{Schwartz Space}}
\subsection{\faen{توزیع‌های معتدل}{Tempered Distributions}}

\section{\faen{پیچش و همگرایی}{Convolution}}
\allocatepages{25}{\faen{فصل ۱۵، بخش ۴: پیچش}{Ch.\ 15, Sec.\ 4: Convolution}}{sectionPhysics}{sec15.4}{technical}

\section{\faen{انتگرال‌های بیضوی در نظریه کنترل}{Elliptic Integrals in Control Theory}}
\allocatepages{20}{\faen{فصل ۱۵، بخش ۵: بیضوی در کنترل}{Ch.\ 15, Sec.\ 5: Elliptic Control}}{sectionPhysics}{sec15.5}{technical}

\chapter{\faen{معادلات تعادل جمعیت}{Population Balance Equations}}
\label{chap:population_balance}

\section{\faen{دینامیک جمعیت}{Population Dynamics}}
\allocatepages{30}{\faen{فصل ۱۶، بخش ۱: دینامیک}{Ch.\ 16, Sec.\ 1: Dynamics}}{chapterMath}{sec16.1}{technical}

\section{\faen{فشرده‌سازی گشتاور}{Moment Compression}}
\allocatepages{30}{\faen{فصل ۱۶، بخش ۲: فشرده‌سازی}{Ch.\ 16, Sec.\ 2: Compression}}{chapterMath}{sec16.2}{technical}

\subsection{\faen{نسبت ۸۹٫۷:۱}{89.7:1 Ratio}}
\subsection{\faen{شرایط Hausdorff}{Hausdorff Conditions}}

\section{\faen{انتقال گشتاور}{Moment Transport}}
\allocatepages{20}{\faen{فصل ۱۶، بخش ۳: انتقال}{Ch.\ 16, Sec.\ 3: Transport}}{chapterMath}{sec16.3}{technical}

\section{\faen{محقق‌پذیری}{Realizability}}
\allocatepages{20}{\faen{فصل ۱۶، بخش ۴: محقق‌پذیری}{Ch.\ 16, Sec.\ 4: Realizability}}{chapterMath}{sec16.4}{technical}

\chapter{\faen{کمی‌سازی عدم‌قطعیت}{Uncertainty Quantification}}
\label{chap:uncertainty_quantification}

\section{\faen{عدم‌قطعیت Epistemic}{Epistemic Uncertainty}}
\allocatepages{25}{\faen{فصل ۱۷، بخش ۱: Epistemic}{Ch.\ 17, Sec.\ 1: Epistemic}}{chapterMath}{sec17.1}{technical}

\section{\faen{عدم‌قطعیت Aleatoric}{Aleatoric Uncertainty}}
\allocatepages{25}{\faen{فصل ۱۷، بخش ۲: Aleatoric}{Ch.\ 17, Sec.\ 2: Aleatoric}}{chapterMath}{sec17.2}{technical}

\section{\faen{انتشار عدم‌قطعیت}{Uncertainty Propagation}}
\allocatepages{15}{\faen{فصل ۱۷، بخش ۳: انتشار}{Ch.\ 17, Sec.\ 3: Propagation}}{chapterMath}{sec17.3}{technical}

\section{\faen{کالیبراسیون}{Calibration}}
\allocatepages{15}{\faen{فصل ۱۷، بخش ۴: کالیبراسیون}{Ch.\ 17, Sec.\ 4: Calibration}}{chapterMath}{sec17.4}{technical}

\chapter{\faen{پردازش سیگنال}{Signal Processing}}
\label{chap:signal_processing}

\section{\faen{تبدیلات ۲D}{2D Separable Transforms}}
\allocatepages{25}{\faen{فصل ۱۸، بخش ۱: تبدیلات}{Ch.\ 18, Sec.\ 1: Transforms}}{sectionPhysics}{sec18.1}{technical}

\subsection{\lat{FFT}}
\subsection{\lat{DCT}}
\subsection{\faen{قرارداد ابعاد}{Dimension Contracts}}

\section{\faen{فیلترهای بیضوی}{Elliptic Filters}}
\allocatepages{25}{\faen{فصل ۱۸، بخش ۲: فیلترها}{Ch.\ 18, Sec.\ 2: Filters}}{sectionPhysics}{sec18.2}{technical}

\section{\faen{پیچش FFT}{FFT Convolution}}
\allocatepages{15}{\faen{فصل ۱۸، بخش ۳: پیچش}{Ch.\ 18, Sec.\ 3: Convolution}}{sectionPhysics}{sec18.3}{technical}

\subsection{\faen{قاعده Padding}{Padding Rule}}

\section{\faen{ادغام حسگر}{Sensor Fusion}}
\allocatepages{15}{\faen{فصل ۱۸، بخش ۴: ادغام}{Ch.\ 18, Sec.\ 4: Fusion}}{sectionPhysics}{sec18.4}{technical}

\chapter{\faen{یکریختی Moment-PBE}{Moment-PBE Isomorphism}}
\label{chap:moment_pbe_iso}

\section{\faen{مبنای نظری}{Theoretical Basis}}
\allocatepages{20}{\faen{فصل ۱۹، بخش ۱: مبنا}{Ch.\ 19, Sec.\ 1: Basis}}{chapterMath}{sec19.1}{technical}

\section{\faen{یکریختی}{Isomorphism}}
\allocatepages{20}{\faen{فصل ۱۹، بخش ۲: یکریختی}{Ch.\ 19, Sec.\ 2: Isomorphism}}{chapterMath}{sec19.2}{technical}

\section{\faen{کاربردها}{Applications}}
\allocatepages{20}{\faen{فصل ۱۹، بخش ۳: کاربردها}{Ch.\ 19, Sec.\ 3: Applications}}{chapterMath}{sec19.3}{technical}

\chapter{\faen{نظریه کنترل}{Control Theory}}
\label{chap:control_theory}

\section{\lat{MPC}}
\allocatepages{20}{\faen{فصل ۲۰، بخش ۱: MPC}{Ch.\ 20, Sec.\ 1: MPC}}{chapterMath}{sec20.1}{technical}

\section{\faen{یادگیری تقویتی}{Reinforcement Learning}}
\allocatepages{20}{\faen{فصل ۲۰، بخش ۲: RL}{Ch.\ 20, Sec.\ 2: RL}}{chapterMath}{sec20.2}{technical}

\section{\faen{بهینه‌سازی توپولوژی}{Topology Optimization}}
\allocatepages{20}{\faen{فصل ۲۰، بخش ۳: توپولوژی}{Ch.\ 20, Sec.\ 3: Topology}}{chapterMath}{sec20.3}{technical}

\chapter{\faen{ادغام شبکه‌های عصبی}{Neural Network Merging}}
\label{chap:nn_merging}

\section{\faen{قوانین Epistemic}{Epistemic Laws}}
\allocatepages{25}{\faen{فصل ۲۱، بخش ۱: قوانین}{Ch.\ 21, Sec.\ 1: Laws}}{chapterMath}{sec21.1}{technical}

\subsection{\faen{قانون مرجع}{Reference Law}}
\subsection{\faen{قانون اندازه‌گیری}{Measurement Law}}
\subsection{\faen{قانون تقارن}{Symmetry Law}}

\section{\faen{تکنیک‌های ادغام}{Merging Techniques}}
\allocatepages{20}{\faen{فصل ۲۱، بخش ۲: تکنیک‌ها}{Ch.\ 21, Sec.\ 2: Techniques}}{chapterMath}{sec21.2}{technical}

\section{\faen{تشخیص دامنه}{Domain Detection}}
\allocatepages{15}{\faen{فصل ۲۱، بخش ۳: دامنه}{Ch.\ 21, Sec.\ 3: Domain}}{chapterMath}{sec21.3}{technical}

\chapter{\faen{مقاله فنی فارسی}{Technical Persian Paper}}
\label{chap:tech_persian}

\allocatepages{80}{\faen{فصل ۲۲: مقاله فنی فارسی}{Chapter 22: Persian Technical Paper}}{sectionPhysics}{chap22}{technical}

%═══════════════════════════════════════════════════════════════════════════
% PART V: ARCHITECTURE (700 pages) — Chapters 23–29
%═══════════════════════════════════════════════════════════════════════════
\partseparator{partArchitecture}{architecture}

\part{\faen{معماری}{Architecture}}
\label{part:architecture}
\hypertarget{part:architecture}{}

\chapter{\faen{مشخصات ISA معماری PICAPD}{PICAPD ISA Specification}}
\label{chap:picapd_isa}

\section{\faen{نمای کلی معماری}{Architecture Overview}}
\allocatepages{15}{\faen{فصل ۲۳، بخش ۱: نمای کلی}{Ch.\ 23, Sec.\ 1: Overview}}{partArchitecture}{sec23.1}{architecture}

\subsection{\faen{اهداف طراحی}{Design Goals}}
\subsection{\faen{اصول معماری}{Architectural Principles}}

\section{\faen{وضعیت معماری}{Architectural State}}
\allocatepages{30}{\faen{فصل ۲۳، بخش ۲: وضعیت}{Ch.\ 23, Sec.\ 2: State}}{partArchitecture}{sec23.2}{architecture}

\subsection{\faen{فایل‌های ثبات}{Register Files}}
\subsection{\faen{ثبات‌های عمومی}{General Purpose Registers (X)}}
\subsection{\faen{ثبات‌های اعشاری}{Floating-Point Registers (F)}}
\subsection{\faen{ثبات‌های رویداد}{Event Registers (E)}}
\subsection{\faen{ثبات‌های گشتاور}{Moment Registers (M)}}
\subsection{\faen{ثبات‌های زمینه}{Context Registers (C)}}
\subsection{\faen{ثبات‌های AGM}{AGM State Registers (A)}}
\subsection{\faen{ثبات‌های اعتماد}{Trust Score Registers (T)}}

\section{\faen{فرمت‌های دستور}{Instruction Formats}}
\allocatepages{35}{\faen{فصل ۲۳، بخش ۳: فرمت‌ها}{Ch.\ 23, Sec.\ 3: Formats}}{sectionISA}{sec23.3}{architecture}

\subsection{\lat{Type-E (Event Control)}}
\subsection{\lat{Type-M (Moment Operations)}}
\subsection{\lat{Type-A (AGM Operations)}}
\subsection{\lat{Type-C (Context/Text Stream)}}
\subsection{\lat{Type-H (Hierarchy)}}
\subsection{\lat{Type-S (Safety/Constraint)}}
\subsection{\lat{Type-V (Variational Mechanics)}}

\section{\faen{تخصیص Opcode}{Opcode Allocation}}
\allocatepages{25}{\faen{فصل ۲۳، بخش ۴: Opcode}{Ch.\ 23, Sec.\ 4: Opcodes}}{sectionISA}{sec23.4}{architecture}

\subsection{\faen{کنترل رویداد}{Event Control (0000011)}}
\subsection{\faen{عملیات گشتاور}{Moment Operations (0001011)}}
\subsection{\faen{AGM/بیضوی}{AGM/Elliptic (0010011)}}
\subsection{\faen{جریان متن}{Text Stream (0011011)}}
\subsection{\faen{سلسله‌مراتب}{Hierarchy (0100011)}}
\subsection{\faen{ایمنی/قید}{Safety/Constraint (0101011)}}
\subsection{\faen{مکانیک تغییرات}{Variational Mechanics (0111011)}}

\section{\faen{مجموعه دستورات}{Instruction Set}}
\allocatepages{40}{\faen{فصل ۲۳، بخش ۵: دستورات}{Ch.\ 23, Sec.\ 5: Instructions}}{sectionISA}{sec23.5}{architecture}

\subsection{\lat{ESET, ECLEAR, ETEST, EWAIT, EBCAST}}
\subsection{\lat{MOM.CALC, MOM.COMP, MOM.XPORT, MOM.REAL}}
\subsection{\lat{AGM.ITER, ELI.COMP, TXF.PASS}}
\subsection{\lat{CTX.SLICE, CTX.AGG, CTX.SYNTH}}
\subsection{\lat{CONS.CHK, BYZ.CONS, TMR.VOTE, SAFE.ROLL}}
\subsection{\lat{LEVAL, LGRAD\_Q, VERLET, HAMILT}}

\section{\faen{مدل حافظه}{Memory Model}}
\allocatepages{20}{\faen{فصل ۲۳، بخش ۶: حافظه}{Ch.\ 23, Sec.\ 6: Memory}}{sectionISA}{sec23.6}{architecture}

\subsection{\faen{سازگاری انتشار}{Release Consistency}}
\subsection{\faen{رابطه Happens-Before}{Happens-Before Relation}}
\subsection{\faen{همگام‌سازی رویداد}{Event Synchronization}}

\section{\faen{مدل برنامه‌نویسی}{Programming Model}}
\allocatepages{15}{\faen{فصل ۲۳، بخش ۷: برنامه‌نویسی}{Ch.\ 23, Sec.\ 7: Programming}}{sectionISA}{sec23.7}{architecture}

\section{\faen{استثناها و وقفه‌ها}{Exceptions \& Interrupts}}
\allocatepages{20}{\faen{فصل ۲۳، بخش ۸: استثناها}{Ch.\ 23, Sec.\ 8: Exceptions}}{sectionISA}{sec23.8}{architecture}

\section{\faen{دستورات مشتق از ریاضیات IIT}{Instructions Derived from IIT Mathematics}}
\allocatepages{10}{\faen{فصل ۲۳، بخش ۹: دستورات IIT}{Ch.\ 23, Sec.\ 9: IIT Instructions}}{sectionISA}{sec23.9}{architecture}

\chapter{\faen{اصلاح ISA}{ISA Rectification}}
\label{chap:isa_rectification}

\section{\faen{شکاف‌های حیاتی}{Critical Gaps}}
\allocatepages{35}{\faen{فصل ۲۴، بخش ۱: شکاف‌ها}{Ch.\ 24, Sec.\ 1: Gaps}}{chapterHardware}{sec24.1}{architecture}

\subsection{\faen{معکوس معنایی Tolerance}{Tolerance Semantic Inversion}}
\subsection{\faen{عدم تطابق فیلد ESET}{ESET Field Mismatch}}
\subsection{\faen{تضادهای ثبات چندمنبعی}{Multi-Source Register Conflicts}}
\subsection{\faen{قوانین Context Flow}{Context Flow Rules}}

\section{\faen{فرمت‌های نقطه ثابت}{Fixed-Point Formats}}
\allocatepages{30}{\faen{فصل ۲۴، بخش ۲: نقطه ثابت}{Ch.\ 24, Sec.\ 2: Fixed-Point}}{chapterHardware}{sec24.2}{architecture}

\subsection{\faen{معناشناسی Qm.n}{Qm.n Semantics}}
\subsection{\faen{رفتار سرریز}{Overflow Behavior}}
\subsection{\faen{حالت‌های گردکردن}{Rounding Modes}}
\subsection{\faen{فرمت‌های کانونیک}{Canonical Formats}}

\section{\faen{عملیات AGM/بیضوی}{AGM/Elliptic Operations}}
\allocatepages{25}{\faen{فصل ۲۴، بخش ۳: AGM}{Ch.\ 24, Sec.\ 3: AGM}}{chapterHardware}{sec24.3}{architecture}

\subsection{\faen{ارزیابی K(k)}{K(k) Evaluation}}
\subsection{\faen{تکرار AGM}{AGM Iteration}}
\subsection{\faen{غیرفعالی تابع انتقال}{Transfer Function Passivity}}

\section{\faen{مرزهای پلتفرم}{Platform Boundaries}}
\allocatepages{30}{\faen{فصل ۲۴، بخش ۴: پلتفرم}{Ch.\ 24, Sec.\ 4: Platform}}{chapterHardware}{sec24.4}{architecture}

\subsection{\faen{توپولوژی EPU}{EPU Topology}}
\subsection{\faen{پروتکل Byzantine}{Byzantine Protocol}}
\subsection{\faen{فرمت‌های DMA}{DMA Formats}}
\subsection{\lat{Write-Ahead Log (WAL)}}

\chapter{\faen{پروفایل پلتفرم STOP-5}{Platform Profile STOP-5}}
\label{chap:platform_profile}

\section{\faen{قراردادهای داده}{Data Contracts}}
\allocatepages{30}{\faen{فصل ۲۵، بخش ۱: قراردادها}{Ch.\ 25, Sec.\ 1: Contracts}}{chapterHardware}{sec25.1}{architecture}

\subsection{\faen{فرمت‌های Bitvector}{Bitvector Formats}}
\subsection{\faen{پایپلاین ادراک}{Perception Pipeline}}

\section{\faen{فرمت‌های نقطه ثابت خودرو}{Automotive Fixed-Point}}
\allocatepages{25}{\faen{فصل ۲۵، بخش ۲: فرمت‌ها}{Ch.\ 25, Sec.\ 2: Formats}}{chapterHardware}{sec25.2}{architecture}

\subsection{\lat{Q1.15 (Radar Samples)}}
\subsection{\lat{Q16.16 (Position/Velocity)}}
\subsection{\lat{Q10.6, Q8.8, Q7.9, Q3.13}}

\section{\faen{مراحل پایپلاین}{Pipeline Stages}}
\allocatepages{45}{\faen{فصل ۲۵، بخش ۳: مراحل}{Ch.\ 25, Sec.\ 3: Stages}}{chapterHardware}{sec25.3}{architecture}

\subsection{\lat{Stage 1: Raw Reception}}
\subsection{\lat{Stage 2: Preprocessing}}
\subsection{\lat{Stage 3: Feature Extraction}}
\subsection{\lat{Stage 4: Multi-Modal Fusion}}
\subsection{\lat{Stage 5: Temporal Integration}}
\subsection{\lat{Stage 6: Model Reduction}}
\subsection{\lat{Stage 7: Decision/Output}}

\chapter{\faen{معماری Queen Bee}{Queen Bee Architecture}}
\label{chap:queen_bee}

\section{\faen{سلسله‌مراتب Worker-Manager-Queen}{Worker-Manager-Queen Hierarchy}}
\allocatepages{30}{\faen{فصل ۲۶، بخش ۱: سلسله‌مراتب}{Ch.\ 26, Sec.\ 1: Hierarchy}}{chapterHardware}{sec26.1}{architecture}

\subsection{\lat{100 Workers}}
\subsection{\lat{10 Managers}}
\subsection{\lat{1 Queen}}

\section{\faen{فشرده‌سازی گشتاور}{Moment Compression}}
\allocatepages{25}{\faen{فصل ۲۶، بخش ۲: فشرده‌سازی}{Ch.\ 26, Sec.\ 2: Compression}}{chapterHardware}{sec26.2}{architecture}

\section{\faen{تصمیم‌گیری جمعیت}{Population Decision}}
\allocatepages{25}{\faen{فصل ۲۶، بخش ۳: تصمیم}{Ch.\ 26, Sec.\ 3: Decision}}{chapterHardware}{sec26.3}{architecture}

\section{\faen{اجماع مقاوم به خطا با فرم‌های دوخطی}{Fault-Tolerant Consensus Using Bilinear Forms}}
\allocatepages{10}{\faen{فصل ۲۶، بخش ۴: اجماع}{Ch.\ 26, Sec.\ 4: Consensus}}{chapterHardware}{sec26.4}{architecture}

\chapter{\faen{طراحی سیلیکون EPU}{EPU Silicon Design}}
\label{chap:epu_silicon}

\allocatepages{70}{\faen{فصل ۲۷: طراحی سیلیکون}{Chapter 27: Silicon Design}}{chapterHardware}{chap27}{architecture}

\chapter{\faen{تحمل خطای بیزانسی}{Byzantine Fault Tolerance}}
\label{chap:byzantine_ft}

\section{\faen{پروتکل ۳-مرحله‌ای}{3-Phase Protocol}}
\allocatepages{30}{\faen{فصل ۲۸، بخش ۱: پروتکل}{Ch.\ 28, Sec.\ 1: Protocol}}{chapterHardware}{sec28.1}{architecture}

\subsection{\lat{Broadcast Phase}}
\subsection{\lat{Echo Phase}}
\subsection{\lat{Ready Phase}}

\section{\faen{پیچیدگی پیام}{Message Complexity}}
\allocatepages{25}{\faen{فصل ۲۸، بخش ۲: پیچیدگی}{Ch.\ 28, Sec.\ 2: Complexity}}{chapterHardware}{sec28.2}{architecture}

\subsection{\lat{O(n\textsuperscript{2}) per Phase}}
\subsection{\faen{تحمل $\lfloor(n-1)/3\rfloor$}{Tolerates $\lfloor(n-1)/3\rfloor$}}

\section{\faen{مدل حافظه}{Memory Model}}
\allocatepages{15}{\faen{فصل ۲۸، بخش ۳: حافظه}{Ch.\ 28, Sec.\ 3: Memory}}{chapterHardware}{sec28.3}{architecture}

\section{\faen{تضمین‌های زمانی}{Timing Guarantees}}
\allocatepages{10}{\faen{فصل ۲۸، بخش ۴: زمان}{Ch.\ 28, Sec.\ 4: Timing}}{chapterHardware}{sec28.4}{architecture}

\subsection{\lat{<1.5\textmu s Total Latency}}

\chapter{\faen{معماری ایمنی}{Safety Architecture}}
\label{chap:safety_architecture}

\section{\faen{غیرقابل‌رمزگذاری}{Non-Encodable}}
\allocatepages{15}{\faen{فصل ۲۹، بخش ۱: غیرقابل‌رمزگذاری}{Ch.\ 29, Sec.\ 1: Non-Encodable}}{chapterHardware}{sec29.1}{architecture}

\section{\faen{غیرقابل‌مسیریابی}{Non-Routable}}
\allocatepages{15}{\faen{فصل ۲۹، بخش ۲: غیرقابل‌مسیریابی}{Ch.\ 29, Sec.\ 2: Non-Routable}}{chapterHardware}{sec29.2}{architecture}

\section{\faen{غیرقابل‌مجوز}{Non-Authorizable}}
\allocatepages{20}{\faen{فصل ۲۹، بخش ۳: غیرقابل‌مجوز}{Ch.\ 29, Sec.\ 3: Non-Authorizable}}{chapterHardware}{sec29.3}{architecture}

%═══════════════════════════════════════════════════════════════════════════
% PART VI: FINANCIAL & BUSINESS (350 pages) — Chapters 30–35
%═══════════════════════════════════════════════════════════════════════════
\partseparator{partFinancial}{financial}

\part{\faen{مالی و کسب‌وکار}{Financial \& Business}}
\label{part:financial}
\hypertarget{part:financial}{}

\chapter{\faen{مدل کسب‌وکار}{Business Model}}
\label{chap:business_model}
\allocatepages{65}{\faen{فصل ۳۰: مدل کسب‌وکار}{Chapter 30: Business Model}}{partFinancial}{chap30}{financial}

\chapter{\faen{استراتژی تأمین مالی}{Funding Strategy}}
\label{chap:funding_strategy}
\allocatepages{55}{\faen{فصل ۳۱: تأمین مالی}{Chapter 31: Funding}}{chapterBusiness}{chap31}{financial}

\chapter{\faen{مسیرهای IPO}{IPO Pathways}}
\label{chap:ipo_pathways}
\allocatepages{80}{\faen{فصل ۳۲: مسیرهای IPO}{Chapter 32: IPO Paths}}{chapterBusiness}{chap32}{financial}

\chapter{\faen{تحلیل ارزش‌گذاری}{Valuation Analysis}}
\label{chap:valuation}
\allocatepages{55}{\faen{فصل ۳۳: ارزش‌گذاری}{Chapter 33: Valuation}}{chapterBusiness}{chap33}{financial}

\chapter{\faen{مالکیت معنوی}{Intellectual Property}}
\label{chap:ip}
\allocatepages{50}{\faen{فصل ۳۴: مالکیت معنوی}{Chapter 34: IP}}{chapterBusiness}{chap34}{financial}

\chapter{\faen{تأثیر اقتصادی}{Economic Impact}}
\label{chap:economic_impact}
\allocatepages{30}{\faen{فصل ۳۵: تأثیر اقتصادی}{Chapter 35: Economic Impact}}{chapterBusiness}{chap35}{financial}

%═══════════════════════════════════════════════════════════════════════════
% PART VII: VALIDATION & RESULTS (400 pages) — Chapters 36–41
%═══════════════════════════════════════════════════════════════════════════
\partseparator{partValidation}{validation}

\part{\faen{اعتبارسنجی و نتایج}{Validation \& Results}}
\label{part:validation}
\hypertarget{part:validation}{}

\chapter{\faen{تست انطباق ISA}{ISA Conformance Testing}}
\label{chap:isa_conformance}

\section{\faen{تست‌های دستور}{Instruction Tests}}
\allocatepages{30}{\faen{فصل ۳۶، بخش ۱: دستورات}{Ch.\ 36, Sec.\ 1: Instructions}}{partValidation}{sec36.1}{validation}

\section{\faen{تست‌های حافظه}{Memory Model Tests}}
\allocatepages{25}{\faen{فصل ۳۶، بخش ۲: حافظه}{Ch.\ 36, Sec.\ 2: Memory}}{partValidation}{sec36.2}{validation}

\section{\faen{تست‌های Byzantine}{Byzantine Tests}}
\allocatepages{25}{\faen{فصل ۳۶، بخش ۳: Byzantine}{Ch.\ 36, Sec.\ 3: Byzantine}}{partValidation}{sec36.3}{validation}

\section{\faen{بردارهای طلایی}{Golden Vectors}}
\allocatepages{25}{\faen{فصل ۳۶، بخش ۴: بردارها}{Ch.\ 36, Sec.\ 4: Vectors}}{partValidation}{sec36.4}{validation}

\chapter{\faen{مجموعه Benchmark}{Benchmark Suite}}
\label{chap:benchmark_suite}

\section{\lat{Perception Benchmarks}}
\allocatepages{25}{\faen{فصل ۳۷، بخش ۱: Perception}{Ch.\ 37, Sec.\ 1: Perception}}{chapterResults}{sec37.1}{validation}

\section{\faen{Benchmark‌های تصمیم}{Decision Benchmarks}}
\allocatepages{25}{\faen{فصل ۳۷، بخش ۲: تصمیم}{Ch.\ 37, Sec.\ 2: Decision}}{chapterResults}{sec37.2}{validation}

\section{\faen{Benchmark‌های ایمنی}{Safety Benchmarks}}
\allocatepages{30}{\faen{فصل ۳۷، بخش ۳: ایمنی}{Ch.\ 37, Sec.\ 3: Safety}}{chapterResults}{sec37.3}{validation}

\chapter{\faen{تحلیل مقایسه‌ای}{Comparative Analysis}}
\label{chap:comparative_analysis}

\section{\faen{مقایسه با رقبا}{Competitor Comparison}}
\allocatepages{30}{\faen{فصل ۳۸، بخش ۱: رقبا}{Ch.\ 38, Sec.\ 1: Competitors}}{chapterResults}{sec38.1}{validation}

\subsection{\lat{Tesla FSD}}
\subsection{\lat{Waymo Driver}}
\subsection{\lat{XPeng}}

\section{\faen{مقایسه معماری}{Architecture Comparison}}
\allocatepages{20}{\faen{فصل ۳۸، بخش ۲: معماری}{Ch.\ 38, Sec.\ 2: Architecture}}{chapterResults}{sec38.2}{validation}

\section{\faen{تحلیل تانسور}{Tensor Analysis}}
\allocatepages{20}{\faen{فصل ۳۸، بخش ۳: تانسور}{Ch.\ 38, Sec.\ 3: Tensor}}{chapterResults}{sec38.3}{validation}

\chapter{\faen{نتایج تجربی}{Experimental Results}}
\label{chap:experimental_results}

\section{\faen{نتایج سیلیکون}{Silicon Results}}
\allocatepages{35}{\faen{فصل ۳۹، بخش ۱: سیلیکون}{Ch.\ 39, Sec.\ 1: Silicon}}{chapterResults}{sec39.1}{validation}

\section{\faen{نتایج شبیه‌سازی}{Simulation Results}}
\allocatepages{35}{\faen{فصل ۳۹، بخش ۲: شبیه‌سازی}{Ch.\ 39, Sec.\ 2: Simulation}}{chapterResults}{sec39.2}{validation}

\section{\faen{اعتبارسنجی میدانی}{Field Validation}}
\allocatepages{30}{\faen{فصل ۳۹، بخش ۳: میدانی}{Ch.\ 39, Sec.\ 3: Field}}{chapterResults}{sec39.3}{validation}

\section{\faen{اعتبارسنجی نگاشت دوخطی بر داده IIT}{Validation of the Bilinear Form Mapping on IIT Dataset}}
\allocatepages{10}{\faen{فصل ۳۹، بخش ۴: نگاشت IIT}{Ch.\ 39, Sec.\ 4: IIT Mapping}}{chapterResults}{sec39.4}{validation}

\section{\faen{نتایج شبیه‌سازی سخت‌افزاری EPU}{EPU Hardware Emulation Results}}
\allocatepages{5}{\faen{فصل ۳۹، بخش ۵: شبیه‌سازی EPU}{Ch.\ 39, Sec.\ 5: EPU Emulation}}{chapterResults}{sec39.5}{validation}

\chapter{\faen{اعتبارسنجی ایمنی}{Safety Validation}}
\label{chap:safety_validation}

\section{\faen{تطابق ASIL-D}{ASIL-D Compliance}}
\allocatepages{25}{\faen{فصل ۴۰، بخش ۱: ASIL-D}{Ch.\ 40, Sec.\ 1: ASIL-D}}{sectionBenchmark}{sec40.1}{validation}

\section{\faen{تست‌های تزریق خطا}{Fault Injection}}
\allocatepages{25}{\faen{فصل ۴۰، بخش ۲: تزریق خطا}{Ch.\ 40, Sec.\ 2: Fault Injection}}{sectionBenchmark}{sec40.2}{validation}

\chapter{\faen{یکپارچه‌سازی سیستم و اعتبارسنجی استقرار}{System Integration and Deployment Validation}}
\label{chap:system_integration}

\section{\faen{پایلوت یکپارچه‌سازی با OEM}{OEM Integration Pilot}}
\allocatepages{10}{\faen{فصل ۴۱، بخش ۱: OEM}{Ch.\ 41, Sec.\ 1: OEM Pilot}}{chapterResults}{sec41.1}{validation}

\section{\faen{سازگاری با خودرو نرم‌افزارمحور}{Software-Defined Vehicle Compatibility}}
\allocatepages{10}{\faen{فصل ۴۱، بخش ۲: SDV}{Ch.\ 41, Sec.\ 2: SDV}}{chapterResults}{sec41.2}{validation}

\section{\faen{ارزیابی آمادگی استقرار}{Deployment Readiness Assessment}}
\allocatepages{5}{\faen{فصل ۴۱، بخش ۳: آمادگی}{Ch.\ 41, Sec.\ 3: Readiness}}{chapterResults}{sec41.3}{validation}

\section{\faen{بازخورد مشتری و درس‌آموخته‌ها}{Customer Feedback and Lessons Learned}}
\allocatepages{5}{\faen{فصل ۴۱، بخش ۴: بازخورد}{Ch.\ 41, Sec.\ 4: Feedback}}{chapterResults}{sec41.4}{validation}

%═══════════════════════════════════════════════════════════════════════════
% PART VIII: APPENDICES & REFERENCES (250 pages)
%═══════════════════════════════════════════════════════════════════════════
\partseparator{partAppendix}{appendices}

\part{\faen{پیوست‌ها و مراجع}{Appendices \& References}}
\label{part:appendices}
\hypertarget{part:appendices}{}

\chapter{\faen{پیوست الف: کدها}{Appendix A: Code}}
\label{app:code}
\allocatepages{70}{\faen{پیوست الف: فهرست کدها}{Appendix A: Code Listings}}{partAppendix}{appA}{appendices}

\chapter{\faen{پیوست ب: نمودارها}{Appendix B: Diagrams}}
\label{app:diagrams}
\allocatepages{50}{\faen{پیوست ب: نمودارها}{Appendix B: Diagrams}}{chapterReference}{appB}{appendices}

\chapter{\faen{پیوست ج: مشتقات}{Appendix C: Derivations}}
\label{app:derivations}
\allocatepages{45}{\faen{پیوست ج: مشتقات ریاضی}{Appendix C: Math Derivations}}{chapterReference}{appC}{appendices}

\chapter{\faen{پیوست د: داده‌ها}{Appendix D: Raw Data}}
\label{app:raw_data}
\allocatepages{40}{\faen{پیوست د: داده‌های خام}{Appendix D: Raw Data}}{chapterReference}{appD}{appendices}

\chapter{\faen{پیوست ه: مراجع}{Appendix E: References}}
\label{app:references}
\allocatepages{45}{\faen{پیوست ه: مراجع و کتاب‌شناسی}{Appendix E: Bibliography}}{chapterReference}{appE}{appendices}

%═══════════════════════════════════════════════════════════════════════════
% FINAL PAGE
%═══════════════════════════════════════════════════════════════════════════
\clearpage
\thispagestyle{empty}

\begin{center}
\vspace*{3cm}

{\Huge\bfseries\color{partExecutive}\lat{Ghost Autonomy}\par}
\vspace{1cm}

{\LARGE\bfseries\color{headerColor}مستندات فنی جامع\\\lat{Complete Technical Documentation}\par}
\vspace{0.5cm}

{\large\color{chapterStrategy}نسخه ۲٫۰ — پیشرفته با ناوبری بصری\\
\lat{Version 2.0 --- Enhanced with Visual Navigation}\par}

\vspace{2cm}

\navmap{none}

\vspace{2cm}

{\large\bfseries
۳۱۴۰+ صفحه — ۴۵ فصل — ۸ بخش\\
\lat{3,140+ Pages --- 45 Chapters --- 8 Parts}
}

\vspace{1cm}

{\color{gray}\rule{0.6\textwidth}{0.5pt}}

\vspace{0.8cm}

{\large\bfseries\color{headerColor}
ویژگی‌های پیشرفته:
\begin{itemize}[leftmargin=2cm]
\item ناوبری بصری هشت‌ضلعی با برجسته‌سازی بخش فعال
\item دکمه بازگشت به فهرست در هر صفحه
\item فهرست مطالب گرانوله با زیربخش‌ها
\item یکپارچه‌سازی کامل PICAPD ISA
\item اصلاحات معماری Stop 1-5
\item فلسفه طراحی و مفهوم‌سازی (بخش جدید)
\item دارایی‌های استراتژیک و مشارکت‌های جهانی (بخش جدید)
\end{itemize}
}

\vspace{0.8cm}

{\color{gray}\footnotesize
\lat{© 2026 Ghost Autonomy. All Rights Reserved.}\\
تاریخ: فوریه ۲۰۲۶ / \lat{Date: February 2026}
}

\end{center}

\end{document}